\documentclass[11pt]{article}

\usepackage[utf8]{inputenc}
\usepackage{color}
\usepackage{graphicx}
\usepackage{geometry}
\usepackage{hyperref}
\hypersetup{
    colorlinks,
    citecolor=black,
    filecolor=black,
    linkcolor=black,
    urlcolor=black
}

\renewcommand{\contentsname}{Inhaltsverzeichnis}

\begin{document}

\begin{titlepage}
	\begin{center}
		\vspace*{1cm}

		\Huge
		\textbf{Tourney}\\
		Spezifikation

		\vspace{0.5cm}
		\LARGE
		Softwarepraktikum\\
		\Large
		Wintersemester 2014/2015

		\vspace{1.5cm}

		\large
		\textbf{Jonas Auer (2860992)\\
				 Fabian Biester (2859084)\\
				 Jan Tagscherer (2893134)}

		\vfill

		\includegraphics[width=0.4\textwidth]{Logo.png}

		\vspace{1.5cm}

		\Large
		Universität Stuttgart\\
		14.11.2014
	\end{center}
\end{titlepage}

\newpage

\tableofcontents
\newpage

\section{Einleitung}

\subsection{Zweck der Spezifikation}

Dieses Dokument dient dem Zweck, sowohl funktionale als auch qualitative Anforderungen überprüfbar, vollständig und systematisch festzuhalten. Damit ermöglicht es dem Kunden und den Betreuern sicherzustellen, dass die richtige Funktionalität auf korrekte Art und Weise umgesetzt wird. Um dies zu gewährleisten muss die Spezifikation als wichtiger Leitfaden während der Entwicklung und für alle weiteren erstellten Artefakte beachtet werden. Zudem müssen sowohl die Spezifikation als auch weitere betreffende Dokumente aktuell gehalten werden.

\subsection{Leserkreis}

Der Leserkreis dieses Dokuments setzt sich aus den folgenden Personengruppen zusammen:
\begin{itemize}
	\item Das Projektteam:
	\begin{itemize}
		\item Die Entwickler
		\item Die Gutachter während des Reviews
		\item Die Betreuer
	\end{itemize}
	\item Der Auftraggeber
	\item Mit der Wartung und Weiterentwicklung betraute Software-Entwickler
\end{itemize}

\subsection{Einsatzbereich und Ziele}

In diesem Projekt wird die Organisationssoftware \textbf{Tourney} umgesetzt. Diese setzt sich zum Ziel, die Planung und Durchführung von Spielturnieren deutlich zu vereinfachen.

Zuvor fand die Organisation dieser Turniere weitestgehend manuell statt, indem Anmeldungen und Ergebnisse auf Papier oder per Tabellenkalkulation aufgezeichnet wurden. Dies führte zu starken Verzögerungen und deutlich mehr Arbeit im Rahmen eines solchen Turniers.

Mit dem Einsatz von \textbf{Tourney} entfällt unnötiger Mehraufwand, indem Events und deren Turniere modular verwaltet werden. Dazu gehören sowohl das Führen von Datenbe-ständen wie Voranmeldungen, Anmeldungen und Turnierergebnissen als auch das automatische Auswerten von Spielresultaten anhand bearbeitbarer Spielregeln und die Verteilung einiger Daten über Turnierordner oder weitere administrative Arbeitsplätze mit anschließendem Zusammenfügen.

Die Software wird im Rahmen des Softwarepraktikums im Wintersemester 2014/2015 an der Universität Stuttgart und unter Auftrag des Heidelberger Spieleverlags entwickelt.

\subsection{Fachbegriffe und Abkürzungen}

Relevante Fachbegriffe und deren Abkürzungen sind in der Spezifikation \textit{kursiv} markiert und werden im Begriffslexikon geklärt, das diesem Dokument anhängt.

\subsection{Übersicht und Aufbau des Dokuments}

Um die zu entwickelnde Software vollständig zu spezifizieren wird das Dokument in die nachfolgenden Kapitel gegliedert:
\begin{itemize}
	\item[] \textbf{Kapitel 2} gibt einen groben Überblick über grundlegende Funktionen, über die die Software verfügen soll, in welchem Umfeld und von welchen Nutzern sie eingesetzt wird und welche Einschränkungen und Annahmen bei der Entwicklung eingegangen werden.
	\item[] \textbf{Kapitel 3} beschreibt die konkreten Anforderungen, die vom Auftraggeber an die Software gestellt werden. Dabei werden zwischen funktionalen und qualitativen Anforderungen unterschieden.
	\item[] \textbf{Kapitel 4} charakterisiert Nutzer und für die Software relevante Personengruppen.
	\item[] \textbf{Kapitel 5} enthält Skizzen als Prototyp der geplanten Benutzeroberfläche.
	\item[] \textbf{Kapitel 6} umfasst grundlegende Anwendungsfälle, die bei der Benutzung der Software auftreten können und die die Interaktion der Nutzer mit dem Programm umreißen.
	\item[] \textbf{Kapitel 7} umschließt ein Begriffslexikon, indem alle wichtigen Fachbegriffe und Abkürzungen geklärt werden.
\end{itemize}

\section{Allgemeine Beschreibung}

\subsection{Einbettung}

\subsection{Einschränkungen bei der Entwicklung}

\subsection{Grundlegende Produktfunktionen}

\subsection{Benutzermerkmale}

\subsection{Annahmen und Abhängigkeiten}

\section{Spezifische Anforderungen}

\subsection{Funktionale Anforderungen}

\subsubsection{Mengengerüst}

\subsubsection{Leistungsanforderungen}

\subsubsection{Externe Schnittstellen}

\subsubsection{...}

\subsection{Qualitätsanforderungen}

\subsubsection{Bedienbarkeit}

\subsubsection{Erweiterbarkeit}

\subsubsection{Sicherheit}

\subsubsection{Zuverlässigkeit}

\subsubsection{Portabilität}

\subsubsection{...}

\section{Benutzerprofile}

\subsection{Administrator}

\subsection{Turnierordner}

\subsection{Spieler}

\section{Benutzeroberfläche}

\subsection{...}

\section{Anwendungsfälle}

\subsection{Übersicht über die Anwendungsfälle}

\subsection{...}

\section{Anhang}

\subsection{Begriffslexikon}

\subsection{Versionshistorie}

\begin{itemize}
	\item Version 1.0 (05.11.2014)
	\begin{itemize}
		\item Grundstruktur der Spezifikation erstellt
	\end{itemize}
\end{itemize}

\end{document}
